\documentclass[11pt]{article}
\usepackage[a4paper, top=18mm, bottom=18mm, left=20mm, right=20mm]{geometry}
\usepackage[T2A]{fontenc}
\usepackage[utf8]{inputenc}
\usepackage[ukrainian]{babel}
\usepackage{graphicx}
\setlength{\parindent}{0pt}
\setlength{\parskip}{4pt}
\pagestyle{empty}
\begin{document}
%begin
{\large\textbf{test}}

\textbf{Варіант \No\,1}\hfill \textit{ПІБ:}\,\underline{\hspace{60mm}}
\vspace{4mm}

%beginitem
1. Що буде виведено за виконання фрагменту коду?
Скільки 2017-елементних підмножин множини натуральних від 1 до 20172017 містять рівно одне число, що більше 172000
%enditem

%beginitem
2. Що буде виведено за виконання фрагменту коду?
Знайти кількість відношень на n-елементній множині, які нерефлексивні та несиметричні
%enditem

%beginitem
3. Що буде виведено за виконання фрагменту коду?
Обчислити значення суми {$\sum\limits_{k=0}^{n} \left({{2^{2n+2k}}+2}\right){C_n^k}$}
%enditem

%beginitem
4. Що буде виведено за виконання фрагменту коду?
Скількома способами 20172018 однакових куль можна розкласти по 2016 різних пакетах так, щоб кількість пакетів з рівно 2 кулями дорівнювала 2, з рівно 6 кулями дорівнювала 2, а у решті пакетів було не менше ніж по 10 куль?
%enditem

%beginitem
5. Що буде виведено за виконання фрагменту коду?
У папці XX розташовано проект, до якого входить синаксично правильна одиниця трансляції 1.cpp з рядком\par
{\#}include "2.h"
\parФайла 2.h у папці XX нема, але він є в папці Y2. Що треба зробити, щоб одиниця трансляції 1.cpp, могла успішно відкомпілюватися?Переміщувати, копіювати та змінювати файли з кодом заборонено.\par\vspace{1.5cm}
%enditem

%beginitem
6. Що буде виведено за виконання фрагменту коду?
Проект складається з од���ї одиниці трансляції 1.cpp, яка починається таким фрагментом\par
long f(){ return myFactorial();}\par
Функція myFactorial(); означена в 2.cpp.\par
Що треба зробити для того, щоб 1.cpp успішно відкомпілювався?\par\vspace{1.5cm}
Що треба зробити для того, щоб за проектом зібрався виконуваний код?\par

%enditem

%beginitem
7. Що буде виведено за виконання фрагменту коду?
Скiльки змiнних типу int* тут оголошено: int *a, &e, **m, *c, &t, **d;
%enditem

%beginitem
8. Що буде виведено за виконання фрагменту коду?
char a[1]; cout<<sizeof(a)/sizeof(char);
%enditem

%beginitem
9. Що буде виведено за виконання фрагменту коду?
char a[]="abcde0"; cout<<sizeof(a)<<','<<a<<endl;
   a[0]=0; cout<<sizeof(a)<<','<<a<<'*';
%enditem

%beginitem
10. Що буде виведено за виконання фрагменту коду?
Задано тип вузла списку struct Node {int n; Node *prev, *next;};
У списку зберігаються послідовні цілі числа від -42 до 43. Вказівник p встановлено на вузол, в якому зберігається число 9.
Чому дорівнює значення p->prev->prev->prev->prev->prev->n ?
%enditem

%beginitem
11. Що буде виведено за виконання фрагменту коду?
В умовах попередньої задачі було виконано p->prev=p->next->prev;
Чому дорівнює значення p->prev->prev->prev->prev->n ?
%enditem

%beginitem
12. Що буде виведено за виконання фрагменту коду?
if (0 < 0) cout<<"a"; else cout<<"a";
%enditem

%beginitem
13. Що буде виведено за виконання фрагменту коду?
int f(int &x, int &y) { x=1; y+=1; return x;}
 int main(){
  int a=1, b=1; a=f(a,a); cout<<a<<":"<<b<<endl;
  {int a=3, b=3; cout<<((a<3) && ((b+=1)<3));cout<<a<<":"<<b<<endl;}
  {inta=5; cout<<a<<endl;}
  cout<<a<<":"<<b<<endl; return 0;}
%enditem

%beginitem
14. Що буде виведено за виконання фрагменту коду?
(Задача за частиною Пашка А.О. Наводиться в авторській редакції.)\par {\hspace {0.5 cm}}³дбудеться деяка подія $W$ чи ні залежить від двох факторів $A$ та $X$. За результатами статистичного експерименту (100 раз для кожного значення фактора $A$,  200 раз для кожного значення фактора $X$) подія $W$ відбулася:\par {\hspace {0.5 cm}}$(n+17)$ раз для значення $A(a)$, $(n+21)$ раз для значення $A(b)$, $(n+50)$ раз для значення $A(c)$, $(n+57)$ раз для значення $A(d)$, та відповідно $(2n+100)$ раз для значення $X(x)$, $(4n+17)$ раз для значення $X(y)$, $(n+150)$ раз для значення $X(z)$.\par
{\hspace {0.5 cm}}Знайти ймовірність того що подія $W$ не відбулася (умовна ймовірність) при значеннях факторів $A(b)$, $X(y)$. Фактори є незалежними, $n$ – номер цього екзаменаційного білета.\par
%enditem

%beginitem
15. Що буде виведено за виконання фрагменту коду?
try:
    res = 3/3*3*3
    if isinstance(res, bool):
        print(f'bool;{res}')
    elif isinstance(res, int):
        print(f'int;{res}')
    elif isinstance(res, float):
        print(f'float;{res:.2f}')
    else:
        print('???')
except:
    print('ERR')

%enditem

%beginitem
16. Що буде виведено за виконання фрагменту коду?
try:
    res = -1**True+-2
    if isinstance(res, bool):
        print(f'bool;{res}')
    elif isinstance(res, int):
        print(f'int;{res}')
    elif isinstance(res, float):
        print(f'float;{res:.2f}')
    else:
        print('???')
except:
    print('ERR')

%enditem

%beginitem
17. Що буде виведено за виконання фрагменту коду?
try:
    res = not2<6 and notTrue==7 and 2.0!=6.0
    if isinstance(res, bool):
        print(f'bool;{res}')
    elif isinstance(res, int):
        print(f'int;{res}')
    elif isinstance(res, float):
        print(f'float;{res:.2f}')
    else:
        print('???')
except:
    print('ERR')

%enditem

%beginitem
18. Що буде виведено за виконання фрагменту коду?
try:
    res = 1<"1.0"
    if isinstance(res, bool):
        print(f'bool;{res}')
    elif isinstance(res, int):
        print(f'int;{res}')
    elif isinstance(res, float):
        print(f'float;{res:.2f}')
    else:
        print('???')
except:
    print('ERR')

%enditem

%beginitem
19. Що буде виведено за виконання фрагменту коду?
def f(n):
    print('f', end='')
    return n

try:
    print(f(9) and f(-4))
except:
    print('ERR')

%enditem

%beginitem
20. Що буде виведено за виконання фрагменту коду?
a = 0
b = 4
c = 5

def f(a):
global c    a = 1
    b += 3
    c = 5
    return a + b + c

a = 2
b = 6
c = 3

try:
    print(f(a), a, b, c, sep=';')
except:
    print('ERR', a, b, c, sep=';')

%enditem

%beginitem
21. Що буде виведено за виконання фрагменту коду?
a = 0
b = 4
c = 5

def f(a):
global c    a = 1
    b = 3
    c = 5
    return a + b + c

a = 2
b = 6
c = 3

try:
    print(f(a), a, b, c, sep=';')
except:
    print('ERR', a, b, c, sep=';')

%enditem

%beginitem
22. Що буде виведено за виконання фрагменту коду?
try:
    i = 0
    while i<5:
        i += 1
    print(i, end=';')
except:
    print('ERR')

%enditem

%beginitem
23. Що буде виведено за виконання фрагменту коду?
try:
    i = 5
    while i>0:
        i += -3
    print(i, end=';')
except:
    print('ERR')

%enditem

%beginitem
24. Що буде виведено за виконання фрагменту коду?
try:
    for a in range(14, 14, -1):
        if a<3:
            continue
        print(a, end=';')
        if a<3:
            break
    else:
        print('end',end=';')
    print(a)
except:
    print('ERR')

%enditem

%beginitem
25. Що буде виведено за виконання фрагменту коду?
try:
    print(0, end=';')
    print(4/1, end=';')
    print(5, end=';')
except Exception:
    print(2, end=';')
except TypeError:
    print(6, end=';')
except:
    print('ERR', end=';')
else:
    print(3, end=';')
finally:
    print(8)

%enditem

%beginitem
26. Що буде виведено за виконання фрагменту коду?
def f(arg, n):
    n = 0
    tmp = arg
    tmp[-9:-9] = [12]
    return len(tmp)>3

try:
    a = [10,11,12,13,14]
    n = 4
    f(a, n)
    print(n, end=';')
    for el in a:
        print(el, end=';')
except:
    print('ERR')

%enditem

%beginitem
27. Що буде виведено за виконання фрагменту коду?
x,y,z=0,4,5;  x,y,z,z=x,6,z,y;  print(x,y,z)
%enditem

%beginitem
28. Що буде виведено за виконання фрагменту коду?
3/3*3*3
%enditem

%beginitem
29. Що буде виведено за виконання фрагменту коду?
-1**True+-2
%enditem

%beginitem
30. Що буде виведено за виконання фрагменту коду?
def f(a):
      u=38
      if a<-5: return 0
      if a==0: u=4
      else: return 5
      return u
   print(f(9))
%enditem

%beginitem
31. Що буде виведено за виконання фрагменту коду?
def f(a,b):
      c=38
      if a<-5: return 0
      if b==0: c=4
      else: return 5
      return c
   print(f(9,9))
%enditem

%beginitem
32. Що буде виведено за виконання фрагменту коду?
not2<6 and notTrue==7 and 2.0!=6.0
%enditem

%beginitem
33. Що буде виведено за виконання фрагменту коду?
1<"1.0"
%enditem

%beginitem
34. Що буде виведено за виконання фрагменту коду?
def f(n): print("f", end=""); return n
   print(f(9) and f(-4))
%enditem

%beginitem
35. Що буде виведено за виконання фрагменту коду?
a,b,c=0,4,5
    def f(a):
      global c      a=1
      b+=3
      c=5
      return a+b+c
    a,b,c=2,6,3
    print(f(a),a,b,c)
%enditem

%beginitem
36. Що буде виведено за виконання фрагменту коду?
a,b,c=0,4,5
    def f(a):
      global c      a=1
      b=3
      c=5
      return a+b+c
    a,b,c=2,6,3
    print(f(a),a,b,c)
%enditem

%beginitem
37. Що буде виведено за виконання фрагменту коду?
i=0
   while i<5: i+=1
   print(i)
%enditem

%beginitem
38. Що буде виведено за виконання фрагменту коду?
i=5
   while i>0: i+=-3
   print(i)
%enditem

%beginitem
39. Що буде виведено за виконання фрагменту коду?
for a in range(14,14,-1):
     if a<3: continue
     print(a,end=' ')
     if a<3: break
   else: print('end',end=' ')
   print(a)
%enditem

%beginitem
40. Що буде виведено за виконання фрагменту коду?
a=('a','b','c','d','e',0,1); print(a[0])
%enditem

%beginitem
41. Що буде виведено за виконання фрагменту коду?
a='0123456789'; print(a[14:14:-1])
%enditem

%beginitem
42. Що буде виведено за виконання фрагменту коду?
a=[10,11,12,13,14]; a.remove(8);print(a)
%enditem

%beginitem
43. Що буде виведено за виконання фрагменту коду?
a=[10,11,12,13,14]; a[-9:-9]=[12]; print(a)
%enditem

%beginitem
44. Що буде виведено за виконання фрагменту коду?
Записати ВИРАЗ, значенням якого є список, що складається з квадратівцілих чисел від 10 до 2018 (включно), що не діляться на 2, записаних в оберненому порядку.
%enditem

%beginitem
45. Що буде виведено за виконання фрагменту коду?
Записати ВИРАЗ, значенням якого є список, що складається з цілих чисел від 10 до 2018 (включно), причому спочатку за зростанням записано всі, що не діляться на 2, а потім решту за зростанням.
%enditem

%beginitem
46. Що буде виведено за виконання фрагменту коду?

\begin{center}
	\adjustimage{max size={0.8\linewidth}{0.8\paperheight}}
	{../BinTree/trees_exp/28.png}
\end{center}
Указати послідовність міток вершин дерева, зображеного на рис., що утвориться в результаті обходу дерева в прямому порядку. ̳тки записувати через пробіл.\par Обхід дерева починається з кореня (на рис. це верхня вершина).
%answer:28 прямому
%enditem

%beginitem
47. Що буде виведено за виконання фрагменту коду?
code
a = 0
h = 1

class A:
    a = 2
    
    def __init__(self):
global a        self.a = 3
        a = 4
        A.a = 5

obj = A()
print(a, h, A.a, obj.a, sep=' ')
code

%enditem

%beginitem
48. Що буде виведено за виконання фрагменту коду?
code
_a = 0

class A:
    _a = 1
    
    def __init__(self):
        self._a = 2
        _a = 3
        A._a = 4

obj = A()
try:
    print(_a, end=' ')
    print(obj._a, end=' ')
    print(A._a, end=' ')
    print(obj._a, end=' ')
    print(obj._A__a, end=' ')
except:
    print('ERR')
code

%enditem

%beginitem
49. Що буде виведено за виконання фрагменту коду?
Змінні next та prev вузла списку позначають наступний та попередній за ним вузол відповідно. Починаючи з голови у список записано послідовні цілі числа від -42 до 43. Нехай змінна p позначає вузол, що зберігає значення 9.
Яке значення зберігається у вузлі p.prev.prev.prev.prev.prev ?
%enditem

%beginitem
50. Що буде виведено за виконання фрагменту коду?
Записати ВИРАЗ, значенням якого є множина, що складається з квадратівцілих чисел від 10 до 2018 (включно), що не діляться на 2.
%enditem

%beginitem
51. Що буде виведено за виконання фрагменту коду?
Записати ВИРАЗ, значенням якого є множина, що складається з цілих чисел від 10 до 2018 (включно), що не діляться на 2.
%enditem

%beginitem
52. Що буде виведено за виконання фрагменту коду?
Записати ВИРАЗ, значенням якого є словник, ключами якого є цілі числа від 10 до 2018 (включно), що не діляться на 2, а ключам відповідають їх квадрати 
%enditem

%beginitem
53. Що буде виведено за виконання фрагменту коду?
a,b,c=0,4,5
   def f(a):
     global c     a=1
     b+=3
     c=5
     return a+b+c
   a,b,c=2,6,3
   print(f(a),a,b,c)
%enditem

%beginitem
54. Що буде виведено за виконання фрагменту коду?
a,b,c=0,4,5
   def f(a):
     global c     a=1
     b=3
     c=5
     return a+b+c
   a,b,c=2,6,3
   print(f(a),a,b,c)
%enditem

%beginitem
55. Що буде виведено за виконання фрагменту коду?
try:
    res = 3/3*3*3
    if isinstance(res, bool):
        print(f'bool;{res}')
    elif isinstance(res, int):
        print(f'int;{res}')
    elif isinstance(res, float):
        print(f'float;{res:.2f}')
    else:
        print('???')
except:
    print('Z')

%enditem

%beginitem
56. Що буде виведено за виконання фрагменту коду?
try:
    res = -1**True+-2
    if isinstance(res, bool):
        print(f'bool;{res}')
    elif isinstance(res, int):
        print(f'int;{res}')
    elif isinstance(res, float):
        print(f'float;{res:.2f}')
    else:
        print('???')
except:
    print('Z')

%enditem

%beginitem
57. Що буде виведено за виконання фрагменту коду?
try:
    res = not2<6 and notTrue==7 and 2.0!=6.0
    if isinstance(res, bool):
        print(f'bool;{res}')
    elif isinstance(res, int):
        print(f'int;{res}')
    elif isinstance(res, float):
        print(f'float;{res:.2f}')
    else:
        print('???')
except:
    print('Z')

%enditem

%beginitem
58. Що буде виведено за виконання фрагменту коду?
try:
    res = 1<"1.0"
    if isinstance(res, bool):
        print(f'bool;{res}')
    elif isinstance(res, int):
        print(f'int;{res}')
    elif isinstance(res, float):
        print(f'float;{res:.2f}')
    else:
        print('???')
except:
    print('Z')

%enditem

%beginitem
59. Що буде виведено за виконання фрагменту коду?
def f(n):
    print('f', end=';')
    return n

try:
    print(f(9) and f(-4))
except:
    print('Z')

%enditem

%beginitem
60. Що буде виведено за виконання фрагменту коду?
a = 0
b = 4
c = 5

def f(a):
global c    a = 1
    b += 3
    c = 5
    return a + b + c

a = 2
b = 6
c = 3

try:
    print(f(a), a, b, c, sep=';')
except:
    print('Z', a, b, c, sep=';')

%enditem

%beginitem
61. Що буде виведено за виконання фрагменту коду?
a = 0
b = 4
c = 5

def f(a):
global c    a = 1
    b = 3
    c = 5
    return a + b + c

a = 2
b = 6
c = 3

try:
    print(f(a), a, b, c, sep=';')
except:
    print('Z', a, b, c, sep=';')

%enditem

%beginitem
62. Що буде виведено за виконання фрагменту коду?
try:
    i = 0
    while i<5:
        i += 1
    print(i, end=';')
except:
    print('Z')

%enditem

%beginitem
63. Що буде виведено за виконання фрагменту коду?
try:
    i = 5
    while i>0:
        i += -3
    print(i, end=';')
except:
    print('Z')

%enditem

%beginitem
64. Що буде виведено за виконання фрагменту коду?
try:
    for a in range(14, 14, -1):
        if a<3:
            continue
        print(a, end=';')
        if a<3:
            break
    else:
        print('end',end=';')
    print(a)
except:
    print('Z')

%enditem

%beginitem
65. Що буде виведено за виконання фрагменту коду?
try:
    print(0, end=';')
    print(4/1, end=';')
    print(5, end=';')
except Exception:
    print(2, end=';')
except TypeError:
    print(6, end=';')
except:
    print('Z', end=';')
else:
    print(3, end=';')
finally:
    print(8)

%enditem

%beginitem
66. Що буде виведено за виконання фрагменту коду?
def f(arg, n):
    n = 0
    tmp = arg
    tmp[-9:-9] = [12]
    return len(tmp)>3

try:
    a = [10,11,12,13,14]
    n = 4
    f(a, n)
    print(n, end=';')
    for el in a:
        print(el, end=';')
except:
    print('Z')

%enditem

%beginitem
67. Що буде виведено за виконання фрагменту коду?

x,y,z=0,4,5
x,y,z,z=x,6,z,y
print(x,y,z)

%enditem

%beginitem
68. Що буде виведено за виконання фрагменту коду?

a,b,c=6,7,9
def f(a,b,c):
    print(a,b,c,end="")

f(0,4)
print(a,b,c)


%enditem

%beginitem
69. Що буде виведено за виконання фрагменту коду?
3/3*3*3
%enditem

%beginitem
70. Що буде виведено за виконання фрагменту коду?
-1**True+-2
%enditem

%beginitem
71. Що буде виведено за виконання фрагменту коду?

try:
    print(0,end="")
    print(4/1,end="")
    print(5,end="")
except Exception: print(2,end="")
except ZeroDivisionError: print(6,end="")
else: print(3,end="")
finally: print(8,end="")

%enditem

%beginitem
72. Що буде виведено за виконання фрагменту коду?

def f(a):
    u=38
    if a<-5: 
        return 0
    if a==0:
         u=4
    else:
         return 5
    return u

print(f(9))

%enditem

%beginitem
73. Що буде виведено за виконання фрагменту коду?

def f(a,b):
    c=38
    if a<-5:
        return 0
    if b==0:
         c=4
    else: 
        return 5
    return c

print(f(9,9))
%enditem

%beginitem
74. Що буде виведено за виконання фрагменту коду?
not2<6 and notTrue==7 and 2.0!=6.0
%enditem

%beginitem
75. Що буде виведено за виконання фрагменту коду?
1<"1.0"
%enditem

%beginitem
76. Що буде виведено за виконання фрагменту коду?

def f(n):
    print("f", end="");
    return n

print(f(9) and f(-4))

%enditem

%beginitem
77. Що буде виведено за виконання фрагменту коду?

a,b,c=0,4,5
def f(a):
global c    a=1
    b+=3
    c=5
    return a+b+c

a,b,c=2,6,3
print(f(a),a,b,c)

%enditem

%beginitem
78. Що буде виведено за виконання фрагменту коду?

a,b,c=0,4,5
def f(a):
global c    a=1
    b=3
    c=5
    return a+b+c

a,b,c=2,6,3
print(f(a),a,b,c)

%enditem

%beginitem
79. Що буде виведено за виконання фрагменту коду?

print('R', end=' ')
i=0
while i<5:
    i+=1
print(i)


%enditem

%beginitem
80. Що буде виведено за виконання фрагменту коду?

print('R', end=' ')
i=5
while i>0:
    i+=-3
print(i)


%enditem

%beginitem
81. Що буде виведено за виконання фрагменту коду?

print('R', end=' ')
for a in range(14,14,-1):
    if a<3:
        continue
        print(a, end=' ')
    if a<3:
        break
    else:
        print('end', end=' ')
    print(a)

%enditem

%beginitem
82. Що буде виведено за виконання фрагменту коду?
for a in range(14,14,-1):
    if a<3:
        continue
        print(a,end=' ')
    if a<3:
        break
else:
    print('end', end=' ')
print(a, end=' ')

%enditem

%beginitem
83. Що буде виведено за виконання фрагменту коду?

print('R', end=' ')
a=('a','b','c','d','e',0,1)
print(a[0])


%enditem

%beginitem
84. Що буде виведено за виконання фрагменту коду?

print('R', end=' ')
a='0123456789'
print(a[14:14:-1])


%enditem

%beginitem
85. Що буде виведено за виконання фрагменту коду?

print('R', end=' ')
a=[10,11,12,13,14]
a.remove(8)
print(a)


%enditem

%beginitem
86. Що буде виведено за виконання фрагменту коду?

print('R', end=' ')
a=[10,11,12,13,14]
a[-9:-9]=[12]
print(a)


%enditem

%beginitem
87. Що буде виведено за виконання фрагменту коду?
Записати ВИРАЗ, значенням якого є список, що складається з квадратівцілих чисел від 10 до 2018 (включно), що не діляться на 2, записаних в оберненому порядку.
%enditem

%beginitem
88. Що буде виведено за виконання фрагменту коду?
Записати ВИРАЗ, значенням якого є список, що складається з цілих чисел від 10 до 2018 (включно), причому спочатку за зростанням записано всі, що не діляться на 2, а потім решту за зростанням.
%enditem

%beginitem
89. Що буде виведено за виконання фрагменту коду?

print('R', end=' ')
i=0
while i<5:
    i+=1
print(i)


%enditem

%beginitem
90. Що буде виведено за виконання фрагменту коду?

print('R', end=' ')
i=5
while i>0:
    i+=-3
print(i)


%enditem

%beginitem
91. Що буде виведено за виконання фрагменту коду?

print('R', end=' ')
for a in range(14,14,-1):
    if a<3:
        continue
        print(a, end=' ')
    if a<3:
        break
    else:
        print('end', end=' ')
    print(a)

%enditem

%beginitem
92. Що буде виведено за виконання фрагменту коду?
for a in range(14,14,-1):
    if a<3:
        continue
        print(a,end=' ')
    if a<3:
        break
    else:
        print('end',end=' ')
    print(a)

%enditem

%beginitem
93. Що буде виведено за виконання фрагменту коду?

print('R', end=' ')
a=('a','b','c','d','e',0,1)
print(a[0])


%enditem

%beginitem
94. Що буде виведено за виконання фрагменту коду?

print('R', end=' ')
a='0123456789'
print(a[14:14:-1])


%enditem

%beginitem
95. Що буде виведено за виконання фрагменту коду?

print('R', end=' ')
a=[10,11,12,13,14]
a.remove(8)
print(a)


%enditem

%beginitem
96. Що буде виведено за виконання фрагменту коду?

print('R', end=' ')
a=[10,11,12,13,14]
a[-9:-9]=[12]
print(a)


%enditem

%beginitem
97. Що буде виведено за виконання фрагменту коду?
Записати ВИРАЗ, значенням якого є множина, що складається з квадратівцілих чисел від 10 до 2018 (включно), що не діляться на 2.
%enditem

%beginitem
98. Що буде виведено за виконання фрагменту коду?
Записати ВИРАЗ, значенням якого є множина, що складається з цілих чисел від 10 до 2018 (включно), що не діляться на 2.
%enditem

%beginitem
99. Що буде виведено за виконання фрагменту коду?
Записати ВИРАЗ, значенням якого є словник, ключами якого є цілі числа від 10 до 2018 (включно), що не діляться на 2, а ключам відповідають їх квадрати 
%enditem

%beginitem
100. Що буде виведено за виконання фрагменту коду?
a=[10,11,12,13,14]; a.remove(8);print(a)
%enditem

%end
\newpage
\end{document}


